\section{EA Test 6B: dam break past a building, field scale}

Test~6B of the Environment Agency benchmark \cite{NeelzPender2013} is the
IMPACT isolated-building dam break (Soares-Fraz\~ao and Zech
\cite{SoaresFrazaoZech2002}) with every length scaled up by 20. A closed
channel 73~m wide holds 8~m of water behind a dam with a 20~m opening, and
0.4~m downstream of it. The dam is removed at $t = 0$ and the flood strikes
a 16~m~$\times$~8~m building set at 63$^\circ$ to the channel, forming a
hydraulic jump in front of it and a wake behind it
(Figure~\ref{fig:ea6b_setup}). Manning's $n = 0.05$, the specified
resolution is 2~m (here triangles of at most 4~m$^2$, about 55,000 in
all), and the run lasts 30~min.

The geometry comes from the benchmark's DEM, in which walls, dam and
building stand 10~m above the floor. The channel is the DEM's open cells,
the two dam blocks are cut out of the boundary, and the building is a
rectangle fitted to its DEM cells and left as a reflective hole in the
mesh. Raising the floor 10~m over the building's DEM cells instead raises
the peak levels by up to 0.08~m. Refining the mesh to 1~m$^2$ and
0.25~m$^2$ triangles (221,000 and 884,000) changes the peak levels by at
most 0.12~m (at G3) and the levels after the first minutes by under 1~cm.

\begin{figure}[h]
\begin{center}
\includegraphics[width=0.9\textwidth]{setup.png}
\end{center}
\caption{EA Test 6B: the channel, and the reservoir, dam, building and
gauges G1--G6 close up.}
\label{fig:ea6b_setup}
\end{figure}

The report shows the results of nine full shallow-water packages (ANUGA
1.1 among them) at G2, G4, G5 and G6 (its Figures 4.27--4.28). Peak depths
downstream of the gate lie within 15--20\% of each other, a 50~cm wide
range, and TUFLOW and XPSTORM run higher than the rest in the
high-momentum region below the gate. The water levels and speeds of
Figures~\ref{fig:ea6b_stage} and~\ref{fig:ea6b_speed} lie within that
spread. The checks (Figure~\ref{fig:ea6b_checks}) are the volume balance,
the order in which the bore reaches the gauges, the draining of the
reservoir, a bound on the largest speed (the report questioned ANUGA's
peak velocities in 2013), level agreement at 30~min, and a stored
baseline per flow algorithm.

\begin{figure}
\begin{center}
\includegraphics[width=0.95\textwidth]{stage_plot.png}
\end{center}
\caption{EA Test 6B: water level at G1--G6 over the first 10 minutes.}
\label{fig:ea6b_stage}
\end{figure}

\begin{figure}
\begin{center}
\includegraphics[width=0.95\textwidth]{speed_plot.png}
\end{center}
\caption{EA Test 6B: speed at G1--G6 over the first 10 minutes.}
\label{fig:ea6b_speed}
\end{figure}

\begin{figure}
\begin{center}
\includegraphics[width=0.9\textwidth]{checks.png}
\end{center}
\caption{EA Test 6B: the checks and their outcome.}
\label{fig:ea6b_checks}
\end{figure}
